\section{个股、市场与行业的共同价格上下文}
\subsection{共同下跌作为预测对象的一部分}
证券损失常来自多层环境。市场整体下行时，即使公司近期没有新事件，较大损失概率也可能增加；行业经营条件变化时，同行证券可能同时承压。新主目标保留这些实际损失，价格模型应能直接读取共同状态，而不是先在标签里解释掉共同下跌，再只研究剩余异常。

这种安排不要求把每次损失分解成已识别的因果来源。市场、行业与个股价格存在内生关联，行业指数可能包含目标证券，因子也只是统计参照。本文将它们作为预测时点可知的环境，检验是否改善未来四档概率；冻结预测后可以按实际市场和行业收益作压力诊断，但不把共同相关写成已经证明的传染机制。

\reportfigure{三类掩码历史形成128维共同价格上下文}{fig:marketindustry}{\FigMarketIndustry}{批准价格编码设计：个股64维、市场32维、行业32维。来源与真实CPU装配仍需独立验收；未来实际损失不做市场、行业或因子扣除。}

\subsection{三类历史的经济含义和掩码}
个股序列保留原始实际日收益，并形成截至$t$的累计、波动和下行等特征。市场序列计划使用经认证的原始市场收益、FF三因子及过去状态。原始市场收益与$Rmrf$不同，后者是市场溢酬，不能直接称为指数收益，也不把某个流通市值加权口径擅称上证指数或沪深300。既有FF3按$2|A$字段组合筛选，避免同日多组合连接，但该组合的官方市场身份和证券范围尚未确认，日历一致不足以证明它代表全部A股。

真实CPU候选适配审计还发现，所请求61个历史日期均存在DRET市场收益候选冲突，累计计数为111,758条冲突证券行；缺失和非法值为0，不等于已经得到统一市场序列。不能按日期任取一行或任意平均，需要官方EXCHANGE、MARKET或市场指数收益身份与序列词典补证。这是61日子集核查，不是全部2329日期的统计，亦非监督预测结果；目前市场路径保留掩码，正式市场行业特征尚未装配。

行业序列采用预测日已公开且生效的历史归属所对应的门类收益。现有可直接核查的是证监会行业门类，不混写成概念板块、申万细分行业或交易制度板块。制造业等门类较粗，模型能读取的是这种粗粒度环境，不能据此声称覆盖所有细分行业冲击。归属未知时行业输入带缺失状态，仍保留个股与市场路径。

定义最近60市场日的三类输入$X^s_{i,t},X^M_t,X^I_{g_{i,t},t}$及逐字段掩码$M^s,M^M,M^I$，分别编码
\begin{align}
 h^s_{i,t}&=\operatorname{GRU}_{64}([X^s\odot M^s;M^s]),\\
 h^M_t&=\operatorname{GRU}_{32}([X^M\odot M^M;M^M]),\\
 h^I_{g_{i,t},t}&=\operatorname{GRU}_{32}([X^I\odot M^I;M^I]),\\
 c_{i,t}&=[h^s_{i,t};h^M_t;h^I_{g_{i,t},t}]\in\mathbb R^{128}.
 \label{eq:context}
\end{align}
短历史保持市场日位置并用有效长度与掩码表达，不跳过缺失日形成不同跨度。市场状态每日期只编码一次，行业状态每门类日期只编码一次，再供当日个股共享；这既降低成本，也避免同一市场输入在日批次中被不同预处理重复解释。

\subsection{过去累计、波动和相对表现}
对来源确认的市场或行业日收益$R_{a,s}$，过去窗口累计采用
\begin{equation}
 C_a^{(w)}(t)=\prod_{s=t-w+1}^{t}(1+R_{a,s})-1,\quad w\in\{5,20\},
\end{equation}
20日波动与下行量可表示为
\begin{equation}
 V_a^{(20)}(t)=\operatorname{std}(R_{a,t-19:t}),\quad
 D_a^{(20)}(t)=\sqrt{\frac1{20}\sum_{s=t-19}^{t}\min(R_{a,s},0)^2}.
\end{equation}
所有窗口只截至$t$。完整窗口和部分有效窗口的处理需在字段合同中区分：完整累计必须具备全部必需收益，部分统计则记录有效数与适用规则。掩码防止缺失被当作零波动，不用未来值补齐。

行业相对市场表现$B_g^{(5)}=C_g^{(5)}-C_M^{(5)}$是两个累计收益的比较，不是一只可交易资产的收益，也不把每日“行业减市场”直接复利。其用途是描述行业相对环境，不从主标签重复扣掉市场或行业。若缺少经核验的历史成分权重，不能自行构造剔除目标后的加权行业指数；目标包含在行业指标中的局限在关系与行业控制解释中保留。

\reporttable{共同价格输入的来源、状态及使用约束}{tab:publiccontext}{
\begin{tabularx}{\linewidth}{p{2.7cm}YY}\toprule
输入对象&主要信息&使用边界\\\midrule
个股&实际收益、过去累计/波动/下行及有效率&历史只截至$t$，短历史不删除成员\\
市场&认证原始市场收益、FF3、过去状态&$Rmrf$不冒充原始收益；同日口径唯一\\
行业&时点已知门类收益、过去状态&公开日与生效日共同确认，当前归属不回填\\
相对行业表现&同窗口行业累计减市场累计&仅为历史环境特征，不改实际损失标签\\
缺失和交易状态&字段掩码、有效长度、来源状态&晚到输入不未来回填，估值规则另核\\\bottomrule
\end{tabularx}}

\subsection{历史归属的双时点查询}
行业历史需要同时考虑资料何时公开和归属何时生效。在历史起点$u$做预测，应查询当时可知且有效的$g_{i,u}$，不能把$t$时最新行业覆盖整条过去个股记录。为描述$t$时当前所属门类环境，可以读取该门类截至$t$的历史行业指数；它不意味着公司过去始终属于该门类。研究将这两种查询区分为“过去各时点身份”与“当前已知门类的历史环境”，在输入合同中明确。

现有行业日收益历史审计为50,067行，其中19个缺失值；行业历史4,393行，其中227个缺失值。这些是既有源表审计统计，不是新成员覆盖数。缺少早期已生效归属、分类标准或执行状态说明，会使某些证券无法认证行业路径。程序可以返回行业缺失并使用市场与个股上下文，但不能用今天的名称或行业来补全历史。更细行业若能可靠补证，可以作为后续扩展；当前不为追求图复杂度引入未经核验的细粒度类别。

\subsection{共同上下文贯穿所有专家和强对照}
价格专家以$c_i$预测四档；事件专家以$[c_i;e_i]$预测；关系专家以$[c_i;r_i]$预测。公告聚合query也来自$c_i$。因此，辅助专家不是在完全不了解市场环境时给出一个孤立判断，再由门随意叠加。它们各自在共同环境下结合不同增量信息，最终可靠性权重决定采用程度。

同信息树、拼接、普通门控和标准MoE均读取相同市场行业信息、事件结构、正文表示和关系统计。价格上下文基线保留三个GRU，只有融合结构变化的比较共享同一编码器预算。额外设置个股价格基线和去市场、去行业控制，用于测量共同环境自身的作用。删除市场分支时，同时删除显式市场收益、FF3和依赖市场的派生特征；删除行业时关闭行业序列与行业派生量。个股价格自然包含共同冲击，不可能通过合法特征删除使它经济上完全“无市场”，控制只针对显式信息权限。

\subsection{共同冲击下的解释与压力评价}
市场和行业可能同时影响多数成员。这时只选每日Top10\%会人为把名单容量固定，不能表示风险广度。研究使用验证确定的绝对$p_5,p_{10}$阈值，让共振日期可以有超过10\%的证券达到预警线，并报告当日预测平均概率及已知实际损失广度。该设计直接对应实际损失主目标，而不需要把共同下跌事后补成第二项未经训练的任务。

行情分组分为预测时点可知与事后压力两类。前者根据过去趋势、波动和下行状态，分界只由训练期确定，用于研究条件表现。后者在冻结测试预测后，按未来五日实际市场或行业下跌分析共振漏报和聚合校准，不返回选模或阈值阶段。前者可作为输入状态，后者是结果解释，二者不混用。行业组成保持重连还帮助判断关系增量是否主要重复行业环境，但不会将共同相关解释为因果效应。
